% ### Erste Seite der Bolzenverbindungen ###
\centering
\section*{\myul{Bolzenverbindungen}}

% Klebefläche für das erste Register:
\begin{tikzpicture}[remember picture, overlay, x=1mm, y=-1mm]
    \begin{scope}[shift={(current page.north west)}]
        \draw (206,79) -- (202,79) -- (202,57) -- (206,57);
        \node[rotate=90, anchor=center] at (204,68) {Bolzen};
    \end{scope}
\end{tikzpicture}

\begin{tikzpicture}[remember picture, overlay]

% --- ERSTE REIHE: Am Seitenursprung anfangen
\setlength{\rowheight}{47mm}

    % A1) Erste Kachel unten anbauen
    \setlength{\cellwidth}{190mm}
    \Tile{A1}{(X1.south west)}{\cellwidth}{\rowheight}
        {
            \vspace{-\TileSep}
            \begin{minipage}{0.68\cellwidth-\TileSep}
                \centering
                \resizebox{\textwidth}{!}{%
                \begin{tabular}{lcc}
                    \xrowht[()]{2mm}\textbf{Einbaufall} & \textbf{Bemessungs-Stangenkraft} & \textbf{Bemessungs-Biegemoment} \\ \hline
                    \xrowht[()]{11mm}\begin{tabular}[c]{@{}l@{}}Einbaufall 1:\\ Spielpassung in Gabel\\ Spielpassung in Stange\end{tabular} & $F_\text{B} = j_\text{s}~j~F_\text{S}$ & $M_\text{b,B} = \dfrac{F_\text{B}~(t_\text{s}+2\,t_\text{G})}{8}$ \\ \hline
                    \xrowht[()]{11mm}\begin{tabular}[c]{@{}l@{}}Einbaufall 2:\\ Übermaßpassung in Gabel\\ Spielpassung in Stange\end{tabular} & $F_\text{B} = j_\text{s}~j~F_\text{S}$ & $M_\text{b,B} = \dfrac{F_\text{B}~t_\text{s}}{8}$ \\ \hline
                    \xrowht[()]{11mm}\begin{tabular}[c]{@{}l@{}}Einbaufall 1:\\ Spielpassung in Gabel\\ Übermaßpassung in Stange\end{tabular} & $F_\text{B} = j_\text{s}~j~F_\text{S}$ & $M_\text{b,B} = \dfrac{F_\text{B}~t_\text{G}}{4}$
                \end{tabular}
                }
            \end{minipage}%
            \begin{minipage}{0.32\cellwidth-\TileSep}
                \centering
                \scriptsize
                \vspace{-2mm}
                \begin{align*}
                    F_\text{B}   &= \text{Bemessungs-Stangenkraft}       \\
                    j_\text{s}   &= \text{Sicherheitsvielfaches}         \\
                    j            &= \text{Sicherheitsfaktor (Allg. 1,5)} \\
                    M_\text{b,B} &= \text{Bemessungs-Biegemoment}        \\
                    t_\text{S}   &= \text{Stangendicke}                  \\
                    t_\text{G}   &= \text{Gabeldicke}
                \end{align*}
            \end{minipage}
        }{0}{0}{1}{0}

% --- ZWEITE REIHE: Unten anbauen
\setlength{\rowheight}{60mm}

    % X2) Eigene Zeile für die Überschrift
    \Tile{X2}{(A1.south west)}{190mm}{8mm}
        {
            \titlebf{Scherung}
        }{0}{0}{0}{0}

    % A2) Erste Kachel unten anbauen
    \setlength{\cellwidth}{95mm}
    \Tile{A2}{(X2.south west)}{\cellwidth}{\rowheight}
        {
            \titlerm{Roloff / Matek}
        }{0}{1}{1}{0}

    % B2) Zweite Kachel rechts anbauen
    \setlength{\cellwidth}{95mm}
    \Tile{B2}{(A2.north east)}{\cellwidth}{\rowheight}
        {
            \titlerm{Handbuch Strukturberechnung (HSB)}
        }{0}{0}{1}{0}

\end{tikzpicture}